\documentclass[10pt,a4paper]{amsart}
\usepackage[margin=19mm]{geometry}
\usepackage{amsmath,amssymb,mathtools}
\usepackage[scr=rsfs,cal=euler]{mathalfa}
\usepackage{xfrac}
\usepackage[unicode,hidelinks,bookmarks=false]{hyperref}
\newcommand{\ie}{i.e.\ }
\newcommand{\cf}{cf.\ }
\newcommand{\resp}{respectively\ }
\newcommand{\Subsection}{\subsection}
\providecommand{\nobreakdash}{\nobreak\hbox{-}}
\setlength{\emergencystretch}{2em}
\multlinegap0pt
\usepackage{amssymb}
\usepackage{mathtools}
\usepackage{booktabs}
\newtheorem{theorem}{Theorem}
\newtheorem{lemma}{Lemma}
\title{Partition Frequency Moments: Modularity and Congruences}
\author{Hartosh Singh Bal}
\date{Source: arXiv:2602.09766v1 (CC BY 4.0)}
\begin{document}
\maketitle

\noindent\emph{Source and licence.} Partition Frequency Moments: Modularity and Congruences. Version arXiv:2602.09766v1 (CC BY 4.0), distributed by arXiv under the Creative Commons Attribution 4.0 International licence, http://creativecommons.org/licenses/by/4.0/, as recorded in the arXiv licence metadata for this exact version and shown on the abstract page https://arxiv.org/abs/2602.09766v1. Full licence notice: \texttt{licenses/arxiv-2602.09766-CC-BY-4.0.txt}.\par\smallskip

\noindent\textit{Editorial scope.} Counted unit: the article's theorem thm:overpart-zero (source0.tex lines 1061--1076). The article's complete original proof of it is delivered with it (source0.tex lines 1078--1110, one \texttt{proof} environment of 33 lines). No mathematics is written, repaired or re-derived: the statement and the proof are the article's own lines, in the article's own notation, and no retained line is substituted or re-flowed.  Retained uncounted as the source-local context the proof needs: lines 1509--1560.  Nothing was omitted from any retained span, so the package carries no omission record.  No retained line contains a citation, so the wrapper carries no bibliography and no bibliography entry.  No retained line contains a figure environment, an image-inclusion command or drawing code, so no image is delivered and the package declares no asset dependency.  Editorial formatting only, in addition: an AMS \texttt{amsart} preamble that restates the article's own declarations and macros used by the retained lines, namely the environments \texttt{theorem}, \texttt{lemma} on separate counters, together with the standard packages those lines need.  The retained environments print in this document as Theorem 1, Lemma 1 in order of appearance, whereas the article prints the same items with the numbering of its own full document.  The complete native source of the article as distributed in the arXiv e-print for this exact version is bundled unchanged as \texttt{sources/194344/source0.tex}.
\begin{lemma}[Half--integral Sturm bound]\label{lem:sturm-half}
Let $L\ge1$ and let
\[
f(\tau)=\sum_{n\ge0} a_n q^n,\qquad q=e^{2\pi i\tau},
\]
be a holomorphic modular form (holomorphic on $\mathfrak{H}$ and at all cusps)
of half--integral weight
\[
w=\frac{k}{2},\qquad k\in\mathbb{Z}_{\ge1}\ \text{odd},
\]
on $\Gamma_0(4L)$ with some Nebentypus character~$\chi$, and assume that its
Fourier coefficients lie in a ring in which reduction modulo $\ell$ makes sense
(e.g.\ $\mathbb{Z}$ or $\mathbb{Z}_{(\ell)}$).  Let $\ell$ be a prime and suppose that
\[
a_n \equiv 0 \pmod{\ell}\qquad\text{for all }0\le n\le B,
\]
where
\[
B
=\left\lfloor \frac{k}{12}\,
   \bigl[\mathrm{SL}_2(\mathbb{Z}):\Gamma_0(4L)\bigr]\right\rfloor.
\]
Then $a_n\equiv0\pmod{\ell}$ for all $n\ge0$; that is, $f\equiv0\pmod{\ell}$.

\medskip
In particular, in the odd--moment cases $m=2k-1$ we consider the half--integral
weight form
\[
\widetilde{\mathcal M}_{2k-1}(\tau)
:=\frac{E_{2k}(\tau)-1}{C_{2k}}\;\eta(\tau)^{-1}.
\]
This is in fact a \emph{holomorphic} modular form on $\Gamma_0(4)$: although
$\eta(\tau)^{-1}$ has a pole at each cusp, for any cusp $\mathfrak a$ of
$\Gamma_0(4)$ the scaled expansion of $E_{2k}$ has constant term $1$, hence
$(E_{2k}-1)$ vanishes at $\mathfrak a$ and cancels the cusp pole of $\eta^{-1}$.
It has $q$--expansion
\[
\widetilde{\mathcal M}_{2k-1}(\tau)=q^{-1/24}\sum_{n\ge0} M_{2k-1}(n)\,q^n.
\]
Fix an odd prime $\ell$ and a residue class $r\bmod \ell$ with $r\not\equiv 0$.
Then the progression projection
\[
\Pi_{\ell,r}\widetilde{\mathcal M}_{2k-1}(\tau)
:=\sum_{n\ge0} M_{2k-1}(\ell n+r)\,q^n
\]
is holomorphic at $\infty$ (indeed, $\widetilde{\mathcal M}_{2k-1}$ is already
holomorphic at $\infty$) and hence is an honest holomorphic modular form of
weight $(2k-\tfrac12)$ on a congruence subgroup of level dividing $4\ell^2$ with
some nebentypus character.  Therefore the Sturm bound above applies to
$\Pi_{\ell,r}\widetilde{\mathcal M}_{2k-1}$ at the safe level $4\ell^2$ (and a
fortiori at any smaller certified level).
\end{lemma}

\begin{theorem}[Certified overpartition zero--class congruences]\label{thm:overpart-zero}
Let $M^{\overline{\ }}_m(n)$ be the overpartition frequency moments defined
above.  Then the following congruences hold for all $n\ge0$:
\[
M^{\overline{\ }}_{5}(5n)\equiv 0 \pmod{5},\qquad
M^{\overline{\ }}_{7}(7n)\equiv 0 \pmod{7},\qquad
M^{\overline{\ }}_{11}(11n)\equiv 0 \pmod{11},\qquad
M^{\overline{\ }}_{13}(13n)\equiv 0 \pmod{13},
\]
and moreover in the next instances of the same Fermat residue class
$m\equiv 1\pmod{\ell-1}$ one has
\[
M^{\overline{\ }}_{9}(5n)\equiv 0 \pmod{5},\qquad
M^{\overline{\ }}_{13}(7n)\equiv 0 \pmod{7}.
\]
\end{theorem}

\begin{proof}
Fix one of the displayed pairs $(m,\ell)$.  As in the ordinary partition case,
one forms the half--integral weight modular form attached to the moment
generating function and then applies progression projection.  Concretely, the
projection
\[
\Pi_{\ell,0}\Big(\sum_{n\ge0}M^{\overline{\ }}_m(n)q^n\Big)
=\sum_{n\ge0} M^{\overline{\ }}_m(\ell n)\,q^n
\]
is a holomorphic modular form of weight $m+\tfrac12$ on a congruence subgroup of
level dividing $4\ell^2$, with coefficients in $\mathbb{Z}_{(\ell)}$; hence the
half--integral Sturm bound in Lemma~\ref{lem:sturm-half} applies at the safe
level $4\ell^2$.

For each $(m,\ell)$ we verified computationally that the first $B$ coefficients
of the projected series vanish modulo $\ell$, where $B$ is the Sturm bound at
level $4\ell^2$.  Equivalently, we checked
$M^{\overline{\ }}_m(\ell n)\equiv 0\pmod{\ell}$ for all $0\le n\le B$.
The explicit bounds and checked ranges were:
\[
\begin{array}{c|c|c}
(m,\ell) & B \text{ at level } 4\ell^2 & \text{checked up to } \ell B \\ \hline
(5,5)   & 165  & 825 \\
(9,5)   & 285  & 1425 \\
(7,7)   & 420  & 2940 \\
(13,7)  & 756  & 5292 \\
(11,11) & 1518 & 16698 \\
(13,13) & 2457 & 31941.
\end{array}
\]
By Lemma~\ref{lem:sturm-half}, the congruence then holds for all $n\ge0$ in each
case, proving the theorem.
\end{proof}

\end{document}
